\documentclass[11pt,a4paper]{article}

% ---------- Packages ----------
\usepackage[utf8]{inputenc}
\usepackage[T1]{fontenc}
\usepackage{amsmath,amssymb,amsthm}
\usepackage{mathtools}
\usepackage{empheq}
\usepackage{bm}
\usepackage{geometry}
\usepackage{booktabs}
\usepackage{array}
\usepackage{siunitx}
\usepackage{hyperref}
\usepackage{xcolor}
\usepackage{enumitem}

\geometry{margin=1in}
\hypersetup{
    colorlinks=true,
    linkcolor=blue!50!black,
    urlcolor=blue!50!black,
    citecolor=blue!50!black
}

% ---------- Macros ----------
\newcommand{\OmegaT}{\Omega_t}
\newcommand{\Rt}{R_t}
\newcommand{\At}{A_t}
\newcommand{\Tt}{T_t}
\newcommand{\thetat}{\theta_t}
\newcommand{\vi}{v_i}
\newcommand{\viss}{v_{i,\text{ss}}}
\newcommand{\bvi}{\bar{v}_i}
\newcommand{\bT}{\bar{T}}
\newcommand{\btheta}{\bar{\theta}}
\newcommand{\blambda}{\bar{\lambda}}
\newcommand{\bCT}{\bar{C}_T}
\newcommand{\Tcmd}{T_t^{\text{cmd}}}
\newcommand{\CTcmd}{C_T^{\text{cmd}}}
\newcommand{\lcmd}{\lambda^{\text{cmd}}}

\title{\textbf{Tail Rotor Dynamics for an RC Helicopter}\\[6pt]
\large Thrust, Induced Inflow, and Nonlinear Lead Compensation}
\author{}
\date{\today}

\begin{document}
\maketitle

\begin{abstract}
\noindent This note develops a control-oriented model of an RC helicopter tail rotor.
Starting from blade-element momentum theory (BEMT), we derive the thrust-versus-pitch
relation, add dynamic inflow via the Pitt--Peters first-order model, invert the
coupled system to obtain the blade pitch required to track a commanded thrust, and
then extend the inverse to handle the static and dynamic nonlinearities of inflow
softening. The result is a nonlinear lead compensator suitable for a tail-rotor
inner loop. Realistic parameter ranges for RC tail rotors are tabulated.
\end{abstract}

\tableofcontents
\vspace{1em}
\hrule
\vspace{1em}

% =====================================================================
\section*{Nomenclature}
\addcontentsline{toc}{section}{Nomenclature}
% =====================================================================

\noindent The notation used throughout the note is collected below. Overbars
denote trim (equilibrium) values; \(\delta(\cdot)\) denotes small perturbations
about trim; a superscript ``cmd'' denotes a commanded (reference) quantity.
SI units are used unless otherwise stated; angles are in radians except where a
table lists degrees.

\vspace{0.5em}
\noindent\textit{Rotor geometry and operating condition}
\begin{center}
\renewcommand{\arraystretch}{1.12}
\begin{tabular}{@{}l l l@{}}
\toprule
Symbol & Description & Units \\
\midrule
\(\Rt\)         & Tail-rotor radius                                  & \si{\metre}\\
\(r\)           & Radial coordinate of a blade element, \(0\le r\le \Rt\) & \si{\metre}\\
\(c\)           & Blade chord (assumed rectangular blade)             & \si{\metre}\\
\(N_b\)         & Number of blades                                    & --\\
\(\At\)         & Rotor-disc area, \(\At=\pi\Rt^2\)                   & \si{\square\metre}\\
\(\sigma\)      & Rotor solidity, \(\sigma = N_b c/(\pi\Rt)\)         & --\\
\(\OmegaT\)     & Rotor angular speed                                 & \si{\radian\per\second}\\
\(\OmegaT\Rt\)  & Blade tip speed                                     & \si{\metre\per\second}\\
\(\rho\)        & Air density (taken as \(\SI{1.225}{\kilogram\per\cubic\metre}\)) & \si{\kilogram\per\cubic\metre}\\
\(Re\)          & Blade-section Reynolds number                       & --\\
\bottomrule
\end{tabular}
\end{center}

\vspace{0.5em}
\noindent\textit{Blade-element aerodynamics}
\begin{center}
\renewcommand{\arraystretch}{1.12}
\begin{tabular}{@{}l l l@{}}
\toprule
Symbol & Description & Units \\
\midrule
\(\thetat\)     & Blade collective pitch angle (control input)        & \si{\radian}\\
\(U_T\)         & Tangential component of local flow at a blade section & \si{\metre\per\second}\\
\(U_P\)         & Perpendicular (axial) component of local flow         & \si{\metre\per\second}\\
\(\phi\)        & Local inflow angle, \(\phi\approx U_P/U_T\)           & \si{\radian}\\
\(\alpha\)      & Effective angle of attack, \(\alpha=\thetat-\phi\)    & \si{\radian}\\
\(a\)           & Two-dimensional lift-curve slope (\(\approx \SI{5.7}{\per\radian}\)) & \si{\per\radian}\\
\(C_\ell\)      & Section lift coefficient, \(C_\ell = a\,\alpha\)      & --\\
\bottomrule
\end{tabular}
\end{center}

\vspace{0.5em}
\noindent\textit{Inflow and free-stream velocities}
\begin{center}
\renewcommand{\arraystretch}{1.12}
\begin{tabular}{@{}l l l@{}}
\toprule
Symbol & Description & Units \\
\midrule
\(\vi\)         & Uniform induced velocity through the rotor disc (state) & \si{\metre\per\second}\\
\(\viss\)       & Steady-state value of \(\vi\) for the current operating point & \si{\metre\per\second}\\
\(v_h\)         & Hover induced velocity, \(v_h=\sqrt{\Tt/(2\rho\At)}\) & \si{\metre\per\second}\\
\(V_\infty\)    & Axial component of the free-stream velocity through the rotor & \si{\metre\per\second}\\
\(V_\perp\)     & In-plane component of the free-stream velocity         & \si{\metre\per\second}\\
\(\mu_z\)       & Non-dimensional axial inflow, \(\mu_z = V_\infty/(\OmegaT\Rt)\) & --\\
\(\mu\)         & Advance ratio (in-plane), \(\mu = V_\perp/(\OmegaT\Rt)\) & --\\
\bottomrule
\end{tabular}
\end{center}

\vspace{0.5em}
\noindent\textit{Thrust and non-dimensional coefficients}
\begin{center}
\renewcommand{\arraystretch}{1.12}
\begin{tabular}{@{}l l l@{}}
\toprule
Symbol & Description & Units \\
\midrule
\(\Tt\)         & Tail-rotor thrust                                   & \si{\newton}\\
\(C_T\)         & Thrust coefficient, \(C_T = \Tt/[\rho\At(\OmegaT\Rt)^2]\) & --\\
\(\lambda\)     & Total inflow ratio, \(\lambda=(\vi+V_\infty)/(\OmegaT\Rt)\) & --\\
\(\lambda_i\)   & Induced-inflow ratio, \(\lambda_i = \vi/(\OmegaT\Rt)\) & --\\
\(\lambda_h\)   & Hover inflow ratio, \(\lambda_h=\sqrt{C_T/2}\)        & --\\
\(\lambda_0\)   & Uniform component of inflow in three-state Pitt--Peters & --\\
\(\lambda_s,\lambda_c\) & Longitudinal/lateral harmonic inflow components & --\\
\(C_{Mx},C_{My}\) & Roll and pitch aerodynamic-moment coefficients on the disc & --\\
\bottomrule
\end{tabular}
\end{center}

\vspace{0.5em}
\noindent\textit{Dynamic inflow model}
\begin{center}
\renewcommand{\arraystretch}{1.12}
\begin{tabular}{@{}l l l@{}}
\toprule
Symbol & Description & Units \\
\midrule
\(\tau_i\)      & First-order inflow time constant (Pitt--Peters)     & \si{\second}\\
\(0.849\)       & Apparent-mass coefficient, \(8/(3\pi)\), in \(\tau_i\) & --\\
\([M]\)         & Apparent-mass matrix in three-state Pitt--Peters    & --\\
\([L]\)         & Inflow-gain matrix in three-state Pitt--Peters      & --\\
\bottomrule
\end{tabular}
\end{center}

\vspace{0.5em}
\noindent\textit{Trim (equilibrium) quantities}
\begin{center}
\renewcommand{\arraystretch}{1.12}
\begin{tabular}{@{}l l l@{}}
\toprule
Symbol & Description & Units \\
\midrule
\(\bT\)         & Trim thrust (typically hover)                       & \si{\newton}\\
\(\bvi\)        & Trim induced velocity                               & \si{\metre\per\second}\\
\(\btheta\)     & Trim blade pitch                                    & \si{\radian}\\
\(\blambda\)    & Trim inflow ratio                                   & --\\
\(\bCT\)        & Trim thrust coefficient                             & --\\
\(\bar{\tau}_i\) & Trim inflow time constant                          & \si{\second}\\
\bottomrule
\end{tabular}
\end{center}

\vspace{0.5em}
\noindent\textit{Linearisation, plant model and compensator}
\begin{center}
\renewcommand{\arraystretch}{1.12}
\begin{tabular}{@{}l l l@{}}
\toprule
Symbol & Description & Units \\
\midrule
\(\delta(\cdot)\) & Small perturbation of a quantity about its trim value & (varies)\\
\(s\)             & Laplace variable                                    & \si{\per\second}\\
\(t\)             & Time                                                & \si{\second}\\
\(K_1\)           & Pitch-to-thrust gain, eq.~\eqref{eq:K1K2}           & \si{\newton\per\radian}\\
\(K_2\)           & Inflow-to-thrust gain, eq.~\eqref{eq:K1K2}          & \si{\newton\second\per\metre}\\
\(k_v\)           & Linearised sensitivity \(\partial\viss/\partial T\), eq.~\eqref{eq:kv} & \si{\metre\per\second\per\newton}\\
\(K_2 k_v\)       & Dimensionless DC inflow-softening coupling          & --\\
\(\beta\)         & Softening factor, \(\beta = 1+K_2 k_v\)             & --\\
\(S(\lambda)\)    & Self-consistent BEMT softening factor, eq.~\eqref{eq:dCT_dtheta} & --\\
\(x\)             & Internal state of the lead compensator              & \si{\radian}\\
\bottomrule
\end{tabular}
\end{center}

\vspace{0.5em}
\noindent\textit{Commanded (reference) quantities}
\begin{center}
\renewcommand{\arraystretch}{1.12}
\begin{tabular}{@{}l l l@{}}
\toprule
Symbol & Description & Units \\
\midrule
\(\Tcmd\)              & Commanded tail-rotor thrust                   & \si{\newton}\\
\(\CTcmd\)             & Commanded thrust coefficient                  & --\\
\(\lcmd\)              & Inflow ratio at the commanded trim, \(\lcmd=\sqrt{\CTcmd/2}\) & --\\
\(\thetat^{\text{ss,cmd}}\) & Steady-state pitch corresponding to \(\CTcmd\), eq.~\eqref{eq:static_inverse} & \si{\radian}\\
\bottomrule
\end{tabular}
\end{center}

\vspace{0.5em}
\noindent\textit{Acronyms}
\begin{center}
\renewcommand{\arraystretch}{1.12}
\begin{tabular}{@{}l l@{}}
\toprule
BEMT  & Blade Element Momentum Theory \\
LPV   & Linear Parameter-Varying \\
DC    & Direct current (steady-state, low-frequency limit) \\
TF    & Transfer function \\
RC    & Radio-controlled \\
\bottomrule
\end{tabular}
\end{center}

\vspace{1em}
\hrule
\vspace{1em}

% =====================================================================
\section{Introduction}
% =====================================================================

The tail rotor of a single-main-rotor helicopter is, in steady flight, a
near-quasi-steady thrust generator: its job is to balance the anti-torque of
the main rotor and, via the pilot's pedal input, to provide directional control.
On a radio-controlled (RC) model helicopter the same physics applies, but
the rotor operates at much lower Reynolds number, much higher rotational
speed, and with a single collective control. From a flight-control perspective,
the relevant question is how the commanded blade pitch \(\thetat\) maps to the
thrust \(\Tt\) actually produced, in both steady and transient conditions.

Two effects make this map non-trivial. First, the thrust produced at a given
pitch depends on the induced velocity \(\vi\) through the rotor: as the rotor
generates lift, it accelerates air downwards through the disc, which changes
the local angle of attack of every blade element and reduces the effective
incidence. Second, the induced velocity itself takes a finite time to
develop, because the wake has inertia; the relationship between \(\vi\) and
\(\Tt\) is therefore dynamic, not algebraic. The combined effect is a
\emph{lead--lag} response in thrust to a step in pitch: an immediate
transient at the open-loop (no-inflow) gain, followed by partial relaxation
to a lower steady value as inflow builds up.

This note develops a compact, control-oriented model of these effects.
Section~1 derives the blade-element thrust relation, an expression linear
in pitch and in induced velocity. Section~2 closes the loop on \(\vi\)
using momentum theory and the first-order Pitt--Peters dynamic-inflow model.
Section~3 inverts the coupled system to obtain the pitch required to track a
thrust command and derives the linearised lead--lag transfer function.
Section~4 shows how the gain and time constant of this plant vary with the
trim point and quantifies the resulting \emph{inflow-softening factor}.
Section~5 presents three forms of nonlinear lead compensator---an LPV
scheduled lead, an exact static inverse, and a two-point interpolant---and
compares their behaviour. Section~6 tabulates parameter ranges for typical
RC-helicopter classes and works a representative numerical example.

% =====================================================================
\section{Tail Rotor Thrust vs.\ Blade Pitch}
% =====================================================================

For a tail rotor with collective pitch \(\thetat\) (no cyclic), we use
\emph{Blade Element Momentum Theory} (BEMT). BEMT combines two independent
views of the same rotor. The \emph{local} (blade-element) view treats each
spanwise strip of the blade as a two-dimensional aerofoil that produces lift
according to its own angle of attack; integrating across the span and around
the disc gives total thrust as a function of \emph{both} the geometric pitch
\(\thetat\) and the induced velocity \(\vi\) flowing through the disc. The
\emph{global} (momentum) view treats the disc as a whole and uses conservation
of momentum on the air column passing through it to fix \(\vi\) given the
thrust it is producing. The two views share \(\vi\) as the coupling variable
and must be solved together. This section develops the local half; the global
half is taken up in Section~\ref{sec:inflow}.

\subsection{Blade Element Theory}

Consider a single blade element---a thin spanwise strip---at radial station
\(r\) on a rectangular blade of chord \(c\), rotating at angular speed
\(\OmegaT\). Two velocity components matter at this element:
\begin{align}
U_T &= \OmegaT\,r \qquad &&\text{(tangential, in-plane)},\\
U_P &= \vi + V_\infty \qquad &&\text{(perpendicular, through-disc)}.
\end{align}
\(U_T\) is simply the speed the element sees from rotation: zero at the hub,
linear in \(r\), peaking at \(\OmegaT\Rt\) at the tip. \(U_P\) is the
\emph{through-disc} flow---air being drawn or pushed perpendicular to the
disc plane. It has two contributions: the induced velocity \(\vi\) produced
by the rotor itself (always present whenever the rotor is producing thrust),
and an external axial component \(V_\infty\) from translation of the airframe
along the rotor shaft (for a tail rotor on a helicopter in straight flight,
this is sideways translation; in hover, \(V_\infty=0\)). The two notations
\(V_\infty\) and \(\mu_z\OmegaT\Rt\) refer to the same quantity, with \(\mu_z\)
its non-dimensional form.

The element therefore sees an apparent wind whose direction is tilted out of
the disc plane by the \emph{inflow angle}
\begin{equation}
\phi = \tan^{-1}\!\left(\frac{U_P}{U_T}\right) \approx \frac{U_P}{U_T},
\end{equation}
where the small-angle approximation holds away from the root since
\(U_P\ll U_T\) over most of the span (typical values: \(U_P\sim\)~few m/s,
\(U_T\sim\)~tens of m/s).

The blade is set at geometric pitch \(\thetat\) relative to the disc plane,
but because the apparent wind is tilted by \(\phi\), the element's effective
angle of attack is
\begin{equation}
\alpha = \thetat - \phi.
\end{equation}
This is the central mechanism of the entire section: \(\thetat\) is what we
command, but induced inflow \emph{steals incidence} from the blade by
increasing \(\phi\), and only the difference \(\alpha\) generates lift. The
thrust loss that appears in the integrated result is the cumulative effect
of this stolen incidence across the span.

Assuming a linear two-dimensional lift curve, \(C_\ell = a\,\alpha\) with
slope \(a\approx \SI{5.7}{\per\radian}\), the elemental thrust per blade is
\begin{equation}
dT = \tfrac{1}{2}\rho\,(U_T^2+U_P^2)\,c\,C_\ell\,dr
   \approx \tfrac{1}{2}\rho\,(\OmegaT r)^2\,c\,a\,(\thetat-\phi)\,dr,
\end{equation}
where in the second form we have dropped \(U_P^2\) against \(U_T^2\)
(consistent with \(\phi\) small) and projected the section lift onto the disc
normal (\(\cos\phi\approx 1\)). Section drag contributes only at second order
in \(\phi\) and is ignored here.

\subsection{Integrated Thrust (Uniform Inflow)}

Total rotor thrust is obtained by summing the elemental contribution over the
span of one blade and multiplying by the number of blades \(N_b\):
\begin{equation}
\Tt \;=\; N_b\!\int_0^{\Rt} dT
     \;=\; \tfrac{1}{2}\rho\,a\,c\,N_b\,\OmegaT^2\!\int_0^{\Rt}\!
           (\thetat - \phi)\,r^2\,dr.
\end{equation}
Assuming \emph{uniform} induced velocity (\(\vi\) constant across the disc, a
standard BEMT simplification), \(\phi\) varies with \(r\) only through its
denominator: \(\phi(r) = U_P/(\OmegaT r) = (\vi+V_\infty)/(\OmegaT r)\). The
integrand splits into two pieces,
\begin{equation}
\int_0^{\Rt}\!\thetat\,r^2\,dr = \frac{\thetat\Rt^3}{3},
\qquad
\int_0^{\Rt}\!\phi(r)\,r^2\,dr
  = \frac{\vi+V_\infty}{\OmegaT}\,\frac{\Rt^2}{2},
\end{equation}
each weighted toward the tip---the first as \(r^2\) (pitch contribution),
the second as \(r\) (inflow loss). Non-dimensionalising by
\(\rho\At(\OmegaT\Rt)^2\) and introducing the rotor solidity
\(\sigma = N_b c/(\pi\Rt)\) and the inflow ratio
\(\lambda = (\vi+V_\infty)/(\OmegaT\Rt)\) gives the standard BEMT result:
\begin{empheq}[box=\fbox]{equation}
C_T \;=\; \frac{\sigma a}{2}\!\left(\frac{\thetat}{3} - \frac{\lambda}{2}\right).
\label{eq:CT_BEMT}
\end{empheq}
The factor \(\tfrac{1}{3}\) on \(\thetat\) reflects the cubic radial weighting
of the pitch contribution; the factor \(\tfrac{1}{2}\) on \(\lambda\) reflects
the quadratic weighting of the inflow loss. The dimensional thrust is recovered as
\begin{equation}
\Tt = C_T\,\rho\,\At\,(\OmegaT\Rt)^2, \qquad \At = \pi\Rt^2.
\end{equation}

For a small RC tail rotor operating at essentially constant \(\OmegaT\)
(governor-controlled), \eqref{eq:CT_BEMT} becomes linear in its two arguments
\(\thetat\) and \(\vi\):
\begin{equation}
\Tt \approx K_1\,\thetat - K_2\,\vi,
\label{eq:linT}
\end{equation}
with (specialising to hover, \(V_\infty = 0\))
\begin{equation}
K_1 = \frac{\rho\At(\OmegaT\Rt)^2\,\sigma a}{6},
\qquad
K_2 = \frac{\rho\At\,\OmegaT\Rt\,\sigma a}{4}.
\label{eq:K1K2}
\end{equation}

\paragraph{Physical interpretation of \(K_1\) and \(K_2\).}
\(K_1\) is the thrust the rotor would produce per radian of collective if no
induced inflow developed: the ``open-loop'' or ``frozen-wake'' pitch
effectiveness. It is what the rotor delivers in the first instant after a
pitch step, before the wake has had time to respond. \(K_2\) is the rate at
which thrust is given back to the induced wake as that wake builds up; it
quantifies how much thrust each m/s of induced velocity costs us. Both
constants scale with the common factor \(\rho\At\sigma a\), but they differ
in their tip-speed dependence: \(K_1\) carries \((\OmegaT\Rt)^2\) while
\(K_2\) carries only \(\OmegaT\Rt\). Their ratio is purely kinematic,
\begin{equation}
\frac{K_2}{K_1} \;=\; \frac{3}{2\,\OmegaT\Rt},
\end{equation}
independent of blade geometry, density, or lift-curve slope. This ratio sets
the scale of the inflow-softening effect quantified in Section~\ref{sec:linsoft}:
faster-spinning rotors retain a larger fraction of their open-loop thrust
gain.

\paragraph{Assumptions in force.}
The derivation above relies on the following simplifications, each of which is
revisited in Section~\ref{sec:params}:
\begin{itemize}[leftmargin=*,itemsep=2pt,topsep=2pt]
\item Uniform induced velocity \(\vi\) across the disc (no spanwise variation,
      no harmonic components).
\item Small inflow angle, \(\phi\ll 1\), justifying \(\tan\phi\approx\phi\),
      \(\cos\phi\approx 1\), and \(U_P^2\ll U_T^2\).
\item Linear two-dimensional lift curve up to stall, with no Reynolds- or
      Mach-number correction.
\item Rectangular, untwisted blade with constant chord \(c\).
\item No tip loss and no root cutout (integration from \(0\) to \(\Rt\)).
\item Thrust contribution from drag is ignored (\(O(\phi^2)\)).
\item Constant rotor speed \(\OmegaT\) (held by the main-rotor governor).
\end{itemize}
The first assumption is the one momentum theory must subsequently repair, by
supplying the \(\vi\) that closes \eqref{eq:CT_BEMT}.

% =====================================================================
\section{Induced Inflow: Momentum Theory}\label{sec:inflow}
% =====================================================================

The blade-element derivation of Section~1 produced
\(C_T = (\sigma a/2)(\thetat/3 - \lambda/2)\), which contains \emph{two}
unknowns: the pitch \(\thetat\) (which we set) and the inflow ratio
\(\lambda\) (which we do not). To close the system we need a second
relationship between thrust and induced velocity. \emph{Momentum theory}
supplies it by treating the rotor disc as a whole and applying conservation
of mass and momentum to the column of air passing through it.

\subsection{The Actuator Disc and the Mass-Flow Argument}

Replace the physical rotor by an idealised \emph{actuator disc}: an
infinitesimally thin permeable surface of area \(\At\) that imparts a
uniform pressure jump \(\Delta p\) to the air passing through it, without
imparting swirl or losses. Draw a streamtube enclosing the disc, extending
from far upstream (station~1, where the flow is undisturbed) through the
disc (station~2, just below it---or station~3 just above, depending on
orientation) and out to far downstream (station~4, where pressure has
returned to ambient). For a rotor in hover the upstream flow is at rest
and the streamtube contracts as the air accelerates; for a rotor in
axial climb the upstream flow has velocity \(V_\infty\).

Three relationships govern this column of air:

\begin{enumerate}[leftmargin=*,itemsep=2pt,topsep=2pt]
\item \emph{Mass conservation.} The mass flow rate
      \(\dot{m}=\rho\At U\) is the same at every station, where \(U\) is
      the local axial velocity. At the disc, \(U = V_\infty + \vi\), where
      \(\vi\) is the induced velocity \emph{at the disc}.
\item \emph{Momentum.} The thrust on the rotor is the reaction to the
      change in axial momentum of the air, evaluated between far upstream
      and far downstream:
      \(\Tt = \dot{m}\,(w_\infty - V_\infty)\), where \(w_\infty\) is
      the axial velocity in the far wake.
\item \emph{Bernoulli, separately above and below the disc.} Energy is
      conserved along streamlines on each side of the disc, but \emph{not}
      across it (the rotor does work on the fluid). Applying Bernoulli
      from station~1 to just above the disc, and from just below the disc
      to station~4, and using the pressure jump \(\Delta p\) at the disc
      gives, after some algebra, the celebrated \emph{Froude result}:
      \begin{equation}
      w_\infty \;=\; V_\infty + 2\,\vi.
      \label{eq:froude}
      \end{equation}
\end{enumerate}

\noindent
The Froude result is non-obvious and worth dwelling on: \emph{half} the
total velocity change is acquired upstream of the disc, and the other
half downstream. The wake therefore contracts to half the disc area and
moves at twice the induced velocity. This factor of two appears in every
momentum-theory thrust expression and is the reason rotors are
\emph{efficient} thrust generators at low disc loading.

Substituting \(w_\infty - V_\infty = 2\vi\) into the momentum equation,
together with \(\dot{m} = \rho\At(V_\infty+\vi)\), gives the
\emph{axial-flight} form of the momentum balance:
\begin{equation}
\Tt \;=\; 2\,\rho\,\At\,\vi\,(V_\infty+\vi).
\label{eq:momentum_axial}
\end{equation}

\subsection{Edgewise Flow: The Full Momentum Balance}

For a tail rotor mounted on a translating airframe, the free-stream is
not generally aligned with the rotor axis. Decompose the airframe
velocity at the rotor into a component along the rotor axis (axial,
\(V_\infty\), which feeds the streamtube directly) and a component in
the disc plane (edgewise, \(V_\perp\), which sweeps air sideways through
the disc). The resultant velocity that the disc accelerates is the
vector sum of the in-plane sweep \(V_\perp\) and the through-disc flow
\(V_\infty + \vi\):
\begin{equation}
U_{\text{disc}} \;=\; \sqrt{(V_\infty+\vi)^2 + V_\perp^2}.
\end{equation}
Glauert's hypothesis takes the mass-flow rate through the disc to be
\(\dot{m} = \rho\At\,U_{\text{disc}}\), and the axial momentum imparted
to be \(2\vi\) (as in the purely axial case). The momentum equation then
generalises to
\begin{empheq}[box=\fbox]{equation}
\Tt \;=\; 2\,\rho\,\At\,\vi\,\sqrt{(V_\infty+\vi)^2 + V_\perp^2}.
\label{eq:momentum}
\end{empheq}
The square-root factor is just the resultant velocity at the disc; the
factor of two is Froude. In the purely axial case (\(V_\perp=0\)) we
recover \eqref{eq:momentum_axial}, and in the hover limit
(\(V_\infty=V_\perp=0\)) the relation simplifies further to
\begin{equation}
\Tt \;=\; 2\,\rho\,\At\,\vi^2 \qquad\Longrightarrow\qquad
v_h \;\equiv\; \sqrt{\frac{\Tt}{2\,\rho\,\At}}.
\label{eq:vh}
\end{equation}

\paragraph{Reading the hover result.}
Equation~\eqref{eq:vh} is the single most-cited result of rotor
aerodynamics. It says induced velocity is set by \emph{disc loading}
\(\Tt/\At\): rotors with larger disc area need less induced velocity for
the same thrust, and induced power \(\Tt\,v_h \propto \Tt^{3/2}/\sqrt{\At}\)
is lower as a result. A tail rotor (small \(\At\), high \(\Tt/\At\))
therefore has \emph{high} \(\vi\) compared with a main rotor of equal
thrust---a numerical example in Section~\ref{sec:params} gives
\(v_h\sim\SI{9}{\metre\per\second}\) for a typical 600-class tail rotor
producing only 8\,N of thrust.

\subsection{Non-Dimensional Form}

Dividing \eqref{eq:momentum} through by \(\rho\At(\OmegaT\Rt)^2\) converts
it to a relation between non-dimensional groups. With the inflow ratios
\(\lambda_i = \vi/(\OmegaT\Rt)\) and \(\lambda = (V_\infty+\vi)/(\OmegaT\Rt)\),
and the edgewise advance ratio \(\mu = V_\perp/(\OmegaT\Rt)\), the result is
\begin{equation}
C_T \;=\; 2\,\lambda_i\,\sqrt{\mu^2 + \lambda^2}
\qquad\Longleftrightarrow\qquad
\lambda_i \;=\; \frac{C_T}{2\,\sqrt{\mu^2 + \lambda^2}}.
\label{eq:momentum_nondim}
\end{equation}
The square-root \(\sqrt{\mu^2+\lambda^2}\) is the non-dimensional
resultant velocity at the disc: \(\mu\) is the in-plane (edgewise)
component, \(\lambda\) the through-disc component, and their hypotenuse
is what the disc actually accelerates.

\paragraph{Dimensional and non-dimensional forms in parallel.}
From here on we move between the dimensional induced velocity \(\vi\)
and the inflow ratio \(\lambda_i\) freely; they encode the same
physical quantity, related by
\(\lambda_i = \vi/(\OmegaT\Rt)\). The dimensional form is convenient
inside the dynamic-inflow ODE (Section~\ref{ssec:dyninflow}), the
non-dimensional form inside BEMT algebra, and the hover form
\(v_h\) for sizing arguments.

\subsection{Closing the Coupled System}

Equation~\eqref{eq:momentum_nondim} and the BEMT result
\eqref{eq:CT_BEMT} are two relations between the same three quantities
(\(\thetat\), \(C_T\), \(\lambda\)). For a given pitch \(\thetat\) they
can be solved simultaneously for the steady-state inflow. In hover
(\(\mu=0\) and \(V_\infty=0\), so \(\lambda=\lambda_i\)),
substituting \(C_T\) from BEMT into the momentum result gives
\begin{equation}
\lambda_i \;=\; \frac{(\sigma a / 2)(\thetat/3 - \lambda_i/2)}{2\lambda_i}
\qquad\Longrightarrow\qquad
\lambda_i^2 + \frac{\sigma a}{8}\lambda_i - \frac{\sigma a}{12}\thetat \;=\; 0.
\end{equation}
This is a quadratic in \(\lambda_i\) whose positive root is the standard
hover-BEMT inflow:
\begin{empheq}[box=\fbox]{equation}
\lambda_i \;=\; \frac{\sigma a}{16}\!\left[
\sqrt{1+\frac{64\,\thetat}{3\,\sigma a}}\;-\;1
\right]
\quad\text{(hover, }\mu=0\text{)}.
\label{eq:lambda_implicit}
\end{empheq}
For small \(\thetat\) the bracket expands as
\(\tfrac{32\thetat}{3\sigma a} - \mathcal{O}(\thetat^2)\), giving
\(\lambda_i \approx 2\thetat/3\)---a useful order-of-magnitude check.

\subsection{Why Steady Theory Is Not Enough: Wake Inertia}\label{ssec:dyninflow}

Equation~\eqref{eq:lambda_implicit} gives the \emph{steady} inflow for a
given pitch, but it ignores time. In reality the wake is a column of air
with mass, and changing the rotor's thrust requires accelerating that
column---which cannot happen instantaneously. If at time \(t=0\) the
pilot steps the pitch up, the blade-element thrust jumps immediately to
its open-loop value \(K_1\,\delta\thetat\) (the wake hasn't moved yet,
so the local angle of attack is still \(\thetat-\phi\) at the old
\(\phi\)). Over the next few milliseconds the induced velocity builds
toward its new steady value, the local angle of attack drops, and the
thrust relaxes to its lower equilibrium. The rotor exhibits a
\emph{lead--lag} response in thrust because the inflow has \emph{lag}.

For control-design purposes the simplest faithful model treats the
uniform component of the induced velocity as a first-order state:
\begin{empheq}[box=\fbox]{equation}
\tau_i\,\dot{\vi} + \vi \;=\; \viss,
\label{eq:pittpeters}
\end{empheq}
where \(\viss\) is the steady-state induced velocity that solves
\eqref{eq:momentum} for the \emph{current} thrust \(\Tt\), and
\(\tau_i\) is the wake-relaxation time constant. This is the
\emph{one-state Pitt--Peters} dynamic inflow model. Equation
\eqref{eq:pittpeters} is the inflow analogue of an RC low-pass filter:
\(\viss\) is the target, \(\vi\) is the state, and \(\tau_i\) sets how
quickly the state chases the target.

\subsection{The Pitt--Peters Time Constant}

What sets \(\tau_i\)? Two mechanisms contribute. The first is the simple
flow-traversal time: an air parcel takes roughly \(\Rt/v_h\) to cross
the rotor disc, so \(\tau_i\) must scale this way on dimensional
grounds alone. The second, more subtle effect is the
\emph{apparent mass} of the air column: when an impermeable disc is
accelerated through fluid, the fluid in its immediate neighbourhood is
accelerated with it, contributing an extra inertia even in the absence
of a true rotor wake. Lamb's classical potential-flow result gives the
apparent mass of a thin circular disc as \(\rho\,(8/3)\,\Rt^3\), a
fraction \(8/(3\pi)\approx 0.849\) of the mass of a sphere of radius
\(\Rt\). Pitt and Peters identified this same coefficient as governing
the lag of the uniform-inflow component, yielding
\begin{empheq}[box=\fbox]{equation}
\tau_i \;=\; \frac{8/(3\pi)}{4\,\OmegaT\,\lambda_i^{\text{trim}}}
       \;=\; \frac{0.849\,\Rt}{4\,v_h},
\label{eq:tau_i}
\end{empheq}
where the two forms are equivalent via
\(\lambda_i^{\text{trim}} = v_h/(\OmegaT\Rt)\) at trim, and the
substitution \(\OmegaT\,\lambda_i^{\text{trim}}\Rt = v_h\) collapses the
first form into the second. The right-hand form is the easier to
interpret: \(\tau_i\) is essentially the time for an air parcel
travelling at the hover induced velocity to cross a quarter of the disc
radius, scaled by the apparent-mass coefficient. Higher disc loading
\(\Rightarrow\) faster \(v_h\) \(\Rightarrow\) shorter \(\tau_i\): hard-working
rotors have snappier wakes.

\paragraph{Caveat on Pitt--Peters time constants.}
The closed-form \eqref{eq:tau_i} is known to underpredict measured
inflow lag on full-scale rotors by a factor of 2--4. For RC tail rotors
it likely underpredicts somewhat less, but the value should be regarded
as a lower bound; see Section~\ref{sec:params}.

\subsection{The Coupled Model}

Combining the blade-element thrust relation with the dynamic-inflow
state gives the complete tail-rotor model of this note:
\begin{empheq}[box=\fbox]{align}
\Tt &= \rho\At(\OmegaT\Rt)^2\,\frac{\sigma a}{2}\!\left(\frac{\thetat}{3}-\frac{\vi+V_\infty}{2\OmegaT\Rt}\right),\label{eq:T_coupled}\\[4pt]
\dot{\vi} &= \frac{1}{\tau_i}\!\left(\frac{\Tt}{2\rho\At\sqrt{(V_\infty+\vi)^2+V_\perp^2}} - \vi\right).\label{eq:vi_coupled}
\end{empheq}
Equation~\eqref{eq:T_coupled} is BEMT, algebraic in \(\thetat\) and
\(\vi\). Equation~\eqref{eq:vi_coupled} is Pitt--Peters: a first-order
ODE in \(\vi\) driven by the instantaneous thrust through the
momentum-theory steady-state map. The pair is a one-state nonlinear
system with input \(\thetat\) and output \(\Tt\), parameterised by the
flight condition \((V_\infty, V_\perp)\) and the rotor constants. Every
subsequent section is either inverting this pair (Section~3), linearising
it (Section~4), or designing a compensator around its linearisation
(Section~5).

\subsection{Three-State Pitt--Peters (Extension)}

The one-state model captures only the \emph{uniform} component of
induced velocity. In forward flight the inflow becomes non-uniform: the
advancing-side blade sees different inflow from the retreating-side
blade, and a longitudinal gradient develops as the wake skews backward
under \(V_\perp\). Pitt and Peters captured these effects with two
additional states---\(\lambda_s\) (longitudinal inflow harmonic,
tilting the inflow front-to-back) and \(\lambda_c\) (lateral harmonic,
tilting it side-to-side)---driven by the aerodynamic pitching and
rolling moments on the disc:
\begin{equation}
[M]\begin{Bmatrix}\dot{\lambda}_0\\\dot{\lambda}_s\\\dot{\lambda}_c\end{Bmatrix}
+ [L]^{-1}\begin{Bmatrix}\lambda_0\\\lambda_s\\\lambda_c\end{Bmatrix}
= \begin{Bmatrix}C_T\\ C_{Mx}\\ C_{My}\end{Bmatrix},
\end{equation}
with apparent-mass matrix
\([M]=\mathrm{diag}\bigl(\tfrac{8}{3\pi},\tfrac{16}{45\pi},\tfrac{16}{45\pi}\bigr)\)
and gain matrix \([L]\) depending on advance ratio and wake skew.
The harmonic components matter for main rotors in forward flight and
for cyclic-pitch control, but for a tail rotor in near-hover, with no
cyclic and only modest edgewise flow from main-rotor wake, the uniform
component dominates and the one-state model \eqref{eq:pittpeters} is
almost always sufficient. We use it for the remainder of this note.

% =====================================================================
\section{Inverse Problem: Blade Pitch from Commanded Thrust}
% =====================================================================

The coupled model of Section~\ref{sec:inflow}---blade-element thrust
\eqref{eq:T_coupled} and dynamic inflow \eqref{eq:vi_coupled}---maps a
pitch input \(\thetat\) to a thrust output \(\Tt\). For control design
we need to run this map \emph{backwards}: given a commanded thrust
\(\Tcmd(t)\), what pitch \(\thetat(t)\) should we apply? The answer
will sit inside the tail-rotor inner loop as a feedforward block: the
outer loop (yaw control) produces a thrust demand, the inverse block
translates it into pitch, and the servo drives the swashplate.

The difficulty is that the plant has \emph{one input, one output, one
state}---the state being the induced velocity \(\vi\). The blade-element
relation \eqref{eq:T_coupled} is algebraic in \((\thetat, \vi)\), but
\(\vi\) is itself governed by a first-order ODE and so carries history.
There are two routes through this: invert algebraically and propagate the
state alongside (Section~\ref{ssec:inst_inverse}), or linearise about a
trim point and work in the Laplace domain
(Sections~\ref{ssec:lin}--\ref{ssec:compensator}). Both have their place;
the linearised form is what allows us to talk about \emph{bandwidth},
\emph{phase margin}, and \emph{lead compensation} in the chapters that
follow.

\subsection{The Exact Instantaneous Inverse}\label{ssec:inst_inverse}

The blade-element relation~\eqref{eq:linT} is algebraically invertible:
solving for \(\thetat\) gives
\begin{empheq}[box=\fbox]{equation}
\thetat(t) \;=\; \frac{\Tcmd(t) + K_2\bigl(\vi(t)+V_\infty(t)\bigr)}{K_1}.
\label{eq:inv_instant}
\end{empheq}
This is the \emph{exact} inverse at every instant, conditional on knowing
the current \(\vi\) and \(V_\infty\): apply pitch \eqref{eq:inv_instant}
and the blade element will, instantaneously, produce thrust \(\Tcmd\).

Two practical issues prevent this from being the end of the story:
\begin{itemize}[leftmargin=*,itemsep=2pt,topsep=2pt]
\item \(\vi\) is rarely measured directly on an RC platform. It must be
      \emph{estimated}, typically by integrating the Pitt--Peters
      ODE~\eqref{eq:pittpeters} in parallel with the controller, driven
      by the commanded thrust or by a thrust estimate.
\item \(V_\infty(t)\) varies with airframe motion (sideways translation
      for a tail rotor). Even when the rotor parameters are stationary,
      the external flow generally is not, and the inverse needs a
      reasonably current estimate of \(V_\infty\) to be accurate.
\end{itemize}
With these caveats, \eqref{eq:inv_instant} together with the inflow
ODE~\eqref{eq:pittpeters} \emph{is} a complete control law---a nonlinear
dynamic inversion. Many practical tail-rotor controllers stop here. The
remainder of the section develops the linearised form because it exposes
the bandwidth structure of the plant and makes the lead-compensator
design transparent.

\subsection{Linearisation About Hover Trim}\label{ssec:lin}

A linearised plant model lets us bring frequency-domain tools---transfer
functions, poles and zeros, Bode plots---to bear on a controller that
will ultimately run on the full nonlinear plant. The cost is that the
linear model is only accurate near the chosen trim point.

\paragraph{Choosing trim.}
By \emph{trim} we mean an equilibrium operating point of the plant: a
constant pitch \(\btheta\) producing a constant thrust \(\bT\) with a
constant induced velocity \(\bvi\), all satisfying both BEMT and
momentum theory simultaneously. For a tail rotor on a hovering
helicopter the natural choice is the trim that balances main-rotor
anti-torque in still air. We specialise to hover (\(V_\infty=V_\perp=0\))
because the algebra is cleanest and because the tail rotor spends most
of its operating life near this condition; the non-hover case
generalises by replacing \(\bvi\) with the hover-equivalent
\(v_h=\sqrt{\bT/(2\rho\At)}\) and adjusting the trim accordingly.

\paragraph{Perturbation convention.}
Write each variable as a constant trim value plus a small perturbation:
\begin{equation}
\thetat = \btheta + \delta\thetat,\quad
\vi = \bvi + \delta\vi,\quad
\Tt = \bT + \delta T.
\end{equation}
We retain terms linear in the \(\delta\)-quantities and discard products
of two perturbations (assumed small\,\(\times\)\,small). The trim values
themselves satisfy the steady model and cancel out of the perturbation
equations, leaving a system in the \(\delta\)-variables alone.

\paragraph{What needs linearising.}
The blade-element relation \eqref{eq:linT}, \(\Tt = K_1\thetat - K_2\vi\),
is \emph{already linear} in its two arguments at fixed \(\OmegaT\); its
perturbation form is simply
\begin{equation}
\delta T \;=\; K_1\,\delta\thetat \;-\; K_2\,\delta\vi.
\label{eq:bemt_perturbed}
\end{equation}
The nonlinearity is entirely in the momentum relation
\(\bT = 2\rho\At\bvi^{\,2}\), which couples thrust to induced velocity
through a square. Differentiating implicitly,
\begin{equation}
\frac{\partial T}{\partial \vi}\bigg|_{\text{trim}}
   \;=\; 4\,\rho\At\,\bvi,
\qquad\Longrightarrow\qquad
\delta\viss \;=\; \frac{1}{4\rho\At\bvi}\,\delta T
            \;\equiv\; k_v\,\delta T,
\label{eq:kv}
\end{equation}
where \(k_v\) is the linearised sensitivity of steady-state induced
velocity to thrust at the chosen trim:
\begin{equation}
k_v \;\equiv\; \frac{\partial \viss}{\partial T}\bigg|_{\text{trim}}
   \;=\; \frac{1}{4\rho\At\bvi}
   \;=\; \frac{\bvi}{2\bT},
\qquad\text{units: \si{\metre\per\second\per\newton}}.
\end{equation}
The two forms come from the same expression via the hover relation
\(\bT=2\rho\At\bvi^{2}\). Physically: \(k_v\) tells us \emph{how much
extra induced velocity each newton of additional thrust will pull
through the disc, in steady state}. It is large for low-thrust trims
(where \(\bvi\) is small and any new thrust must accelerate a slow wake
from near-rest) and small for high-thrust trims (where \(\bvi\) is
already substantial and incremental thrust matters less proportionally).

\paragraph{Linearised inflow dynamics.}
The Pitt--Peters ODE \eqref{eq:pittpeters} is already linear in \(\vi\)
at fixed \(\tau_i\). Substituting \(\viss = \bvi + k_v\,\delta T\) and
subtracting the trim equation \(\bvi = \bvi\) leaves
\begin{equation}
\tau_i\,\dot{\delta\vi} + \delta\vi \;=\; k_v\,\delta T.
\end{equation}
Taking the Laplace transform with zero initial conditions
(\(\dot{(\cdot)} \to s\,(\cdot)\)):
\begin{equation}
(\tau_i s + 1)\,\delta\vi(s) \;=\; k_v\,\delta T(s).
\label{eq:dvi_laplace}
\end{equation}

\subsection{The Linearised Plant Transfer Function}\label{ssec:plant_tf}

Equations~\eqref{eq:bemt_perturbed} and \eqref{eq:dvi_laplace} are a
two-equation linear system in the three Laplace-domain quantities
\(\delta\thetat(s)\), \(\delta T(s)\), \(\delta\vi(s)\). Eliminate
\(\delta\vi\): from \eqref{eq:dvi_laplace},
\begin{equation}
\delta\vi(s) \;=\; \frac{k_v}{\tau_i s + 1}\,\delta T(s).
\end{equation}
Substitute into \eqref{eq:bemt_perturbed}:
\begin{equation}
\delta T(s) \;=\; K_1\,\delta\thetat(s)
\;-\;
K_2\,\frac{k_v}{\tau_i s + 1}\,\delta T(s).
\end{equation}
Collect \(\delta T\):
\begin{equation}
\delta T(s)\!\left[1 + \frac{K_2 k_v}{\tau_i s + 1}\right]
\;=\; K_1\,\delta\thetat(s),
\end{equation}
and rearrange to standard form,
\begin{empheq}[box=\fbox]{equation}
\frac{\delta T(s)}{\delta\thetat(s)}
 \;=\; \frac{K_1\,(\tau_i s + 1)}{\tau_i s + (1 + K_2 k_v)}.
\label{eq:plant_tf}
\end{empheq}
This is the linearised tail-rotor plant: thrust per radian of pitch
perturbation about hover trim, valid for small \(\delta\thetat\) over
a bandwidth limited by the linearisation.

\subsection{Reading the Plant: Lead--Lag Structure}\label{ssec:lead_lag}

The plant \eqref{eq:plant_tf} has one zero and one pole, both on the
negative real axis:
\begin{equation}
\text{zero at }s=-\frac{1}{\tau_i},
\qquad
\text{pole at }s=-\frac{1+K_2 k_v}{\tau_i}.
\end{equation}
Because \(K_2 k_v > 0\), the pole sits at a higher frequency than the
zero. The high-frequency and DC limits of the gain are
\begin{align}
\lim_{s\to\infty} \frac{\delta T}{\delta\thetat}
 &= K_1
 \qquad\text{(instantaneous, ``open-loop'' gain)},\\
\lim_{s\to 0} \frac{\delta T}{\delta\thetat}
 &= \frac{K_1}{1+K_2 k_v}
 \qquad\text{(steady-state, ``closed-wake'' gain)}.
\end{align}
\emph{The DC gain is lower than the instantaneous gain by a factor
\(1+K_2 k_v\).} A step in pitch produces an immediate thrust spike of
size \(K_1\,\delta\thetat\) (before the wake responds), which then
relaxes over a timescale \(\tau_i/(1+K_2 k_v)\) to the lower steady
value \(K_1\,\delta\thetat/(1+K_2 k_v)\). This is the \emph{lead--lag}
response described qualitatively in the introduction, now in closed form.

\paragraph{Physical reading of \(K_2 k_v\).}
The dimensionless product \(K_2 k_v\) measures how strongly induced
inflow couples back into the thrust. Tracing its factors:
\(K_2\) is how much thrust each m/s of induced velocity steals; \(k_v\)
is how much extra induced velocity each newton of thrust pulls
through the disc. Their product closes the loop and gives the fractional
thrust loss between the instantaneous and steady-state response:
\begin{equation}
\frac{\text{DC gain}}{\text{instantaneous gain}}
   \;=\; \frac{1}{1+K_2 k_v}.
\end{equation}
For typical RC tail rotors at hover trim, \(K_2 k_v\) sits in the range
\(0.4\)--\(1.0\)---meaning the DC gain is anywhere from \(50\%\) to
\(70\%\) of the open-loop gain. The thrust the rotor delivers in steady
state is significantly less than what its blade pitch suggests. This
phenomenon is \emph{inflow softening}, and quantifying it across the
operating envelope is the subject of Section~\ref{sec:linsoft}.

\paragraph{Caveat: naive vs.\ self-consistent \(K_2 k_v\).}
Computing \(K_2\) from \eqref{eq:K1K2} and \(k_v\) from \eqref{eq:kv}
\emph{separately} and multiplying gives an answer that overstates the
coupling, because the two factors are not independent at trim---they
both depend on \(\bvi\). A self-consistent expression, derived from the
coupled BEMT--momentum equilibrium in Section~\ref{sec:linsoft},
replaces the naive product with a simpler closed form. For now we use
\(K_2 k_v\) as a symbolic placeholder and refine its numerical value in
the next chapter.

\subsection{Inverting the Plant: The Lead Compensator}\label{ssec:compensator}

A feedforward compensator that cancels the plant dynamics is, by
definition, the inverse of the plant. Reciprocating \eqref{eq:plant_tf}:
\begin{empheq}[box=\fbox]{equation}
\frac{\delta\thetat(s)}{\delta T(s)}
 \;=\; \frac{\tau_i s + (1 + K_2 k_v)}{K_1\,(\tau_i s + 1)}.
\label{eq:comp_tf}
\end{empheq}
The roles of pole and zero are now swapped: the compensator has its
\emph{zero} where the plant had its \emph{pole}, and vice versa. The
compensator is in the classical sense a \emph{lead network}: its zero
sits at a lower frequency than its pole, so its phase contribution is
positive across the intervening band. Cascading
\eqref{eq:comp_tf}~\(\to\)~\eqref{eq:plant_tf} gives the identity
transfer function, so applying the compensator's output as pitch
produces exactly the commanded thrust at every frequency---at least
within the validity of the linearisation.

\paragraph{Time-domain reading.}
In the time domain, \eqref{eq:comp_tf} responds to a thrust-command step
by first applying an overdriving pitch transient
(\(\delta\thetat \to \delta T/K_1 \cdot (1+K_2k_v)\) immediately---the
high-frequency limit of the compensator), which then decays to the
steady pitch \(\delta T/K_1\) (the low-frequency limit). In words: the
compensator anticipates the rotor's tendency to fall back as inflow
builds, and pre-emptively over-pitches by the softening factor to
compensate. This is the linearised version of the nonlinear lead
introduced in Section~\ref{sec:linsoft} and refined in Section~6.

\paragraph{Limitations of the linear compensator.}
The compensator \eqref{eq:comp_tf} has gain \(K_1\), zero
\(\beta/\tau_i\), and pole \(1/\tau_i\), all of which depend on the trim
point through \(\bvi\) (and hence through \(\bT\)). A single fixed
compensator designed at hover will mis-track at other thrust levels by
amounts quantified in Section~\ref{sec:linsoft}. Section~6 develops
three constructions---an LPV-scheduled lead, an exact nonlinear static
inverse, and a two-point interpolant---that handle this variation
across the operating envelope.

% =====================================================================
\section{Inflow Softening and the Nonlinear Lead}\label{sec:linsoft}
% =====================================================================
\label{sec:soft}

\subsection{What Inflow Softening Means}

The DC gain \(K_1/(1+K_2k_v)\) varies with the trim point because
\(k_v\propto 1/\bvi \propto 1/\sqrt{\bT}\). The factor \((1+K_2k_v)\) is the
\emph{inflow softening factor}---softening is strongest at low thrust.

\subsection{BEMT Derivation of the Softening Factor}

In hover, the coupled equations are
\begin{align}
C_T &= \frac{\sigma a}{2}\!\left(\frac{\thetat}{3}-\frac{\lambda}{2}\right), \tag{BEMT}\\
\lambda &= \sqrt{C_T/2}. \tag{Momentum}
\end{align}
Eliminating \(C_T\):
\begin{equation}
\lambda^2 = \frac{\sigma a}{4}\!\left(\frac{\thetat}{3}-\frac{\lambda}{2}\right),
\end{equation}
with closed form
\begin{equation}
\lambda(\thetat) = \frac{\sigma a}{16}\!\left[\sqrt{1+\frac{64\thetat}{3\sigma a}}-1\right].
\label{eq:lambda_of_theta}
\end{equation}

\paragraph{Trim sensitivity.} Differentiating implicitly:
\begin{equation}
2\lambda\,\frac{d\lambda}{d\thetat} = \frac{\sigma a}{4}\!\left(\frac{1}{3}-\frac{1}{2}\frac{d\lambda}{d\thetat}\right)
\;\;\Longrightarrow\;\;
\frac{d\lambda}{d\thetat} = \frac{\sigma a}{24\lambda + 3\sigma a/2}.
\end{equation}
Then
\begin{empheq}[box=\fbox]{equation}
\frac{dC_T}{d\thetat}\bigg|_{\text{steady}}
= \frac{\sigma a}{6}\cdot\underbrace{\frac{1}{1+\dfrac{\sigma a}{16\lambda}}}_{S(\lambda)}.
\label{eq:dCT_dtheta}
\end{empheq}
The factor \(\sigma a/6\) is the instantaneous (high-frequency) gain;
\(S(\lambda)\) is the softening factor. Comparing with the linearised
form \(S = 1/(1+K_2 k_v)\) from Section~\ref{ssec:lead_lag} identifies
the self-consistent value of the coupling product as
\begin{empheq}[box=\fbox]{equation}
K_2 k_v \big|_{\text{self-consistent}} \;=\; \frac{\sigma a}{16\,\lambda_h},
\qquad \lambda_h = \sqrt{C_T/2}.
\label{eq:K2kv_consistent}
\end{empheq}
This is the expression to use in design: it bakes in the trim
consistency between \(K_2\) and \(k_v\) that a naive separate
calculation would miss.

\subsection{Magnitude Across the Envelope}

For a 600-class tail rotor (\(\sigma=0.138\), \(a=5.7\), hover \(\bCT=0.024\),
\(\blambda=0.110\)):

\begin{center}
\renewcommand{\arraystretch}{1.15}
\begin{tabular}{lcccc}
\toprule
Thrust & \(C_T\) & \(\lambda\) & \(K_2k_v=\sigma a/(16\lambda)\) & \(S(\lambda)\) \\
\midrule
25\% hover  & 0.006 & 0.055 & 0.89 & 0.53 \\
50\% hover  & 0.012 & 0.078 & 0.63 & 0.61 \\
Hover       & 0.024 & 0.110 & 0.45 & 0.69 \\
150\% hover & 0.036 & 0.135 & 0.36 & 0.73 \\
200\% hover & 0.048 & 0.156 & 0.32 & 0.76 \\
\bottomrule
\end{tabular}
\end{center}

The DC gain varies by roughly \(40\%\) across the operating envelope. A
fixed-gain lead compensator will mis-track by this amount at the extremes.

\subsection{Time Constant Varies Too}

The Pitt--Peters time constant
\(\tau_i = 0.849/(4\OmegaT\lambda)\) scales as \(1/\lambda\propto 1/\sqrt{C_T}\).
For the same 600-class rotor:

\begin{center}
\renewcommand{\arraystretch}{1.15}
\begin{tabular}{lc}
\toprule
Thrust & \(\tau_i\) [ms] \\
\midrule
25\% hover  & 5.5 \\
Hover       & 2.7 \\
200\% hover & 1.9 \\
\bottomrule
\end{tabular}
\end{center}

Both gain and time constant of the rotor depend on the trim operating point.
Both must be scheduled.

% =====================================================================
\section{Nonlinear Lead Compensator}
% =====================================================================

\subsection{Approach A: Gain Scheduling on Commanded Thrust (LPV)}

Schedule the lead-compensator coefficients on the current trim \(\lcmd\)
computed from the commanded thrust:
\begin{equation}
\lcmd(t) = \sqrt{\frac{\CTcmd(t)}{2}}
        = \sqrt{\frac{\Tcmd(t)}{2\rho\At(\OmegaT\Rt)^2}}.
\end{equation}
The scheduled compensator is
\begin{empheq}[box=\fbox]{equation}
\begin{aligned}
\tau_i(t) &= \frac{0.849}{4\,\OmegaT\,\lcmd(t)},\\[4pt]
\beta(t) &\equiv 1+K_2k_v = 1 + \frac{\sigma a}{16\,\lcmd(t)},\\[4pt]
\frac{\thetat(s)}{\Tt(s)} &= \frac{\tau_i(t)\,s+\beta(t)}{K_1\bigl(\tau_i(t)\,s+1\bigr)}.
\end{aligned}
\end{empheq}
In continuous state form, with a single internal state \(x\),
\begin{align}
\dot{x} &= -\frac{1}{\tau_i(t)}\,x + \frac{\beta(t)-1}{\tau_i(t)\,K_1}\,\Tcmd,\\
\thetat &= \frac{1}{K_1}\,\Tcmd + x.
\end{align}
The \(1/K_1\) feedforward gives the instantaneous inversion; the state \(x\)
adds the low-frequency correction.

\subsection{Approach B: Exact Static Inversion + Inflow Filter}

Bypass linearisation. The steady-state map \(\thetat\leftrightarrow C_T\)
is exact in closed form: from~\eqref{eq:CT_BEMT} with \(\lambda=\sqrt{C_T/2}\),
\begin{empheq}[box=\fbox]{equation}
\thetat^{\text{ss,cmd}} \;=\; \frac{6\,\CTcmd}{\sigma a} + \frac{3\,\lcmd}{2},
\qquad \lcmd = \sqrt{\CTcmd/2}.
\label{eq:static_inverse}
\end{empheq}
This static inverse, applied directly as feedforward, intentionally over-drives
the rotor for the duration of the inflow lag: the rotor produces a brief thrust
spike that builds inflow up to \(\lcmd\), after which thrust settles to \(\CTcmd\).
This is the nonlinear analogue of the lead compensator, and is structurally
simpler than the LPV form: a single algebraic expression in \(\CTcmd\).

For completeness, an inflow state with scheduled time constant can be
propagated alongside,
\begin{equation}
\dot{\lambda} = \frac{1}{\tau_i(\lcmd)}\bigl(\lcmd-\lambda\bigr),
\end{equation}
to expose the actual inflow as a state---useful for simulation, monitoring, or
hybridising with feedback.

\subsection{Approach C: Two-Point Approximation}

For embedded implementations where LPV is too heavy, interpolate the
compensator zero \(\beta/\tau_i\) and pole \(1/\tau_i\) linearly in
\(\sqrt{\Tcmd}\) between two design points (low and high thrust).

\subsection{Comparison}

\begin{center}
\renewcommand{\arraystretch}{1.2}
\begin{tabular}{p{4cm}p{3cm}p{3cm}p{4cm}}
\toprule
Compensator & Initial \(\thetat\) & Settled \(\thetat\) & Notes \\
\midrule
None (\(1/K_1\) only) & correct & undershoots ~31\% & thrust slow to build \\
Fixed lead at hover & correct & DC tracks at hover & overshoots at high thrust \\
LPV lead (A) & correct & DC tracks everywhere & needs schedule \\
Exact static inverse (B) & exceeds \(\theta^{\text{ss,cmd}}\) by factor \(\beta(\lcmd)\) (intentional lead) & DC tracks exactly & cleanest, single expression\\
\bottomrule
\end{tabular}
\end{center}

The static-inverse approach (B) is generally preferred: the nonlinearity falls
out of the BEMT algebra without a schedule, and the dynamic correction is a
single first-order state. The LPV form (A) is more rigorous for control analysis
but rarely justifies its complexity in a tail-rotor inner loop.

\subsection{Scheduling Variable Choice}

Should the schedule key on commanded thrust \(\Tcmd\) or on the actual
(estimated) thrust \(\Tt\)? Both have drawbacks. Scheduling on the command
assumes the rotor is tracking; scheduling on the measured thrust is
self-consistent but susceptible to noise. A practical compromise is to schedule
on a low-pass filtered command, with cutoff near \(1/\bar{\tau}_i\)---this is
``scheduling on slow variables'' in LPV terminology and is sound provided the
scheduling-variable bandwidth lies below the loop bandwidth.

% =====================================================================
\section{Parameter Ranges for RC Tail Rotors}\label{sec:params}
% =====================================================================

\subsection{Geometric and Operating Parameters}

\begin{center}
\renewcommand{\arraystretch}{1.15}
\begin{tabular}{llcccc}
\toprule
Symbol & Quantity & Micro & 450 & 550--700 & 800+ \\
\midrule
\(\Rt\) & Tail radius [m]      & 0.03--0.05  & 0.06--0.08  & 0.10--0.13  & 0.14--0.17 \\
\(c\)    & Blade chord [m]       & 0.008--0.012 & 0.015--0.020 & 0.022--0.028 & 0.028--0.035 \\
\(N_b\)  & Number of blades       & 2 & 2 & 2 (occ.\ 4) & 2 or 4 \\
\(\OmegaT\) & Rotor speed [\si{\radian\per\second}] & 800--1500 & 600--1000 & 500--800 & 400--700 \\
\(\OmegaT\Rt\) & Tip speed [m/s] & 40--60 & 50--75 & 70--100 & 80--110 \\
\(\sigma\) & Solidity              & 0.10--0.16 & 0.12--0.18 & 0.12--0.18 & 0.12--0.20 \\
\(a\)    & Lift-curve slope [1/rad] & 5.5--5.75 & 5.5--5.75 & 5.5--5.75 & 5.5--5.75 \\
\bottomrule
\end{tabular}
\end{center}

\noindent The lift-curve slope sits near \(2\pi\approx 6.28\) in theory,
reduced to roughly 5.7 by finite-aspect-ratio and tip-loss effects, and falls
further at very low Reynolds (\(Re<\num{50000}\)).

\subsection{Aerodynamic Trim Values}

\begin{center}
\renewcommand{\arraystretch}{1.15}
\begin{tabular}{llcccc}
\toprule
Symbol & Quantity & Micro & 450 & 550--700 & 800+ \\
\midrule
\(\bT\)      & Hover thrust [N]              & 0.3--1 & 1--4 & 4--12 & 12--30 \\
\(\bvi\)     & Hover induced vel.\ [m/s]     & 4--8 & 5--10 & 7--13 & 9--15 \\
\(\btheta\)  & Trim pitch [deg]              & 6--10 & 6--10 & 6--10 & 6--10 \\
\(\tau_i\)   & Inflow time constant [ms]     & 4--10 & 8--20 & 15--35 & 25--60 \\
\bottomrule
\end{tabular}
\end{center}

\subsection{Derived Coefficients}

Using \(\rho = \SI{1.225}{\kilogram\per\cubic\metre}\):

\begin{center}
\renewcommand{\arraystretch}{1.15}
\begin{tabular}{llcccc}
\toprule
Symbol & Quantity & Micro & 450 & 550--700 & 800+ \\
\midrule
\(\At\)      & Disc area [\si{\square\metre}]   & 0.003--0.008 & 0.011--0.020 & 0.031--0.053 & 0.062--0.091 \\
\(K_1\)      & Pitch gain [N/rad]              & 3--10  & 10--40  & 40--150 & 150--400 \\
\(K_2\)      & Inflow gain [N\,s/m]            & 0.3--1.0 & 0.8--2.5 & 2.5--8 & 8--20 \\
\(k_v\)      & Inflow sensitivity [m/(s\,N)]   & 0.5--2.0 & 0.3--1.0 & 0.2--0.6 & 0.15--0.4 \\
\(K_2 k_v\)  & DC softening factor            & 0.4--1.0 & 0.4--1.0 & 0.4--1.0 & 0.4--1.0 \\
\bottomrule
\end{tabular}
\end{center}

\noindent The dimensionless product \(K_2 k_v\) is roughly the same across
classes---a physical-similarity result, since both factors scale with similar
combinations of \(\rho\At\), \(\OmegaT\Rt\), and \(\bvi\).

\subsection{Worked Example: 600-class}

\begin{align*}
\Rt &= 0.115~\text{m}, & c &= 0.025~\text{m}, & N_b &= 2,\\
\OmegaT &= 700~\text{rad/s}, & \OmegaT\Rt &= 80.5~\text{m/s},\\
\sigma &= \tfrac{2\cdot 0.025}{\pi\cdot 0.115} = 0.138, & a &= 5.7,\\
\bT &= 8~\text{N}, & \At &= 0.0415~\text{m}^2.
\end{align*}
Hence
\begin{align*}
\bvi &= \sqrt{8/(2\cdot 1.225\cdot 0.0415)} = 8.9~\text{m/s},\\
\tau_i &= 0.849\cdot 0.115/(4\cdot 8.9) \approx 2.7~\text{ms (Pitt--Peters)},\\
K_1 &= \tfrac{1.225\cdot 0.0415\cdot 80.5^2\cdot 0.138\cdot 5.7}{6} = 43~\text{N/rad},\\
K_2 &= \tfrac{1.225\cdot 0.0415\cdot 80.5\cdot 0.138\cdot 5.7}{4} = 8.1~\text{N\,s/m},\\
k_v &= 1/(4\cdot 1.225\cdot 0.0415\cdot 8.9) = 0.55~\text{m/(s\,N)},\\
\lambda_h &= \sqrt{0.024/2} = 0.110,\\
K_2k_v\big|_\text{consistent} &= \frac{\sigma a}{16\lambda_h} = \frac{0.787}{1.76} = 0.45.
\end{align*}

\subsection{Caveats}

Pitt--Peters time constants are systematically low compared to flight-test
data on full-scale rotors (factor of 2--4 underestimate); RC practitioners
often use \(\tau_i\approx 0.1\)--\(0.3\) rotor revolutions empirically.
For design, the Pitt--Peters value should therefore be treated as a lower
bound on the inflow lag, and either inflated by a fitted multiplier or
verified against bench data before being used to set compensator
coefficients.
Blade lift-curve slope drops noticeably at low Reynolds numbers
(\(Re\sim\num{30000}\)--\num{80000}) typical of RC tail-rotor tips,
especially with thin symmetric sections. Solidity calculations assume
rectangular blades; tapered or paddle blades need an equivalent solidity
weighted by \(r^2 c(r)\). Finally, the tail rotor sits in the main-rotor
wake, so \(V_\perp\) in~\eqref{eq:momentum} is dominated by main-rotor
inflow---the main source of tail-rotor effectiveness loss in turns.

% =====================================================================
\section{Summary}
% =====================================================================

The tail-rotor inner loop reduces to:

\begin{enumerate}[leftmargin=*]
\item A near-linear blade-element thrust law,
      \(\Tt = K_1\thetat - K_2\vi\), with constants from~\eqref{eq:K1K2}.
\item A first-order inflow state,
      \(\tau_i\dot{\vi}+\vi=\viss(\Tt)\), with \(\tau_i\) from~\eqref{eq:tau_i}.
\item A coupled steady-state map \(\thetat\leftrightarrow C_T\) with closed
      form~\eqref{eq:lambda_of_theta}.
\item A self-consistent DC softening factor \(\sigma a/(16\lambda_h)\) that
      replaces the naive product \(K_2 k_v\).
\item A nonlinear feedforward inverse---either the LPV lead
      \eqref{eq:comp_tf} scheduled on \(\lcmd\), or the exact static
      inverse~\eqref{eq:static_inverse}---that compensates for inflow
      softening across the full thrust envelope.
\end{enumerate}

\end{document}
